% ============================================================================
% Fichier  : partie4/P04_C12_donnees_fraude.tex
% Titre    : Analyse de Données et Détection de la Fraude
% Auteur   : MINKA MI NGUIDJOÏ Thierry Emmanuel
% Version  : 0.1 -- novembre 2026
% ============================================================================
% Sources : INF3226 C06 (fraude) et C27 (audit continu) ; COSO 2013
% principe 8 ; guide COSO-ACFE (dossier Normes) ; Cressey (1953) ; Benford
% (1938) ; Nigrini (2012).
% Faits CAPNO (LAB-06) : 4 312 agents payés en décembre 2025 ; 37 absents
% du fichier du personnel, net mensuel cumulé 9 842 500 FCFA ; 12 comptes
% bancaires partagés par 27 agents ; 9 agents de plus de 60 ans au
% 31/12/2025 encore payés.
% Faits MINÉQUIR (LAB-07) : 23 bons de commande entre 4 800 000 et
% 4 999 999 FCFA en 2025, adressés à 3 fournisseurs, sous le seuil de
% 5 000 000 FCFA (seuil fictif de mise en concurrence de la DAF) ; chiffre 4
% en premier chiffre : 15,8 % des montants des bons (attendu 9,7 %).
% ============================================================================

\chapter{Analyse de Données et Détection de la Fraude}
\label{chap:donnees-fraude}
\chaptermark{Données et fraude}

\epigraph{%
    « Les données ne mentent pas~; elles attendent qu'on leur pose la
    bonne question. »%
}{Adage des auditeurs de données}

\niveauChapitre{Expert}{Chapitres~\ref{chap:preuve}
et~\ref{chap:applications}.}

\begin{boxobjectifs}
\begin{itemize}
    \item Expliquer le risque de fraude au sens du COSO (principe~8) et le
          triangle de la fraude.
    \item Conduire sur un fichier de paie les tests d'existence, d'unicité
          et de cohérence.
    \item Repérer le fractionnement des commandes sous un seuil.
    \item Appliquer la loi de Benford à une série de montants et en
          interpréter l'écart avec prudence.
    \item Passer d'un indice à un constat, sans confondre anomalie et
          fraude.
\end{itemize}
\end{boxobjectifs}

\bigskip

% ============================================================================
\section*{Pourquoi ce chapitre}
% ============================================================================

Les chapitres précédents ont testé des contrôles. Celui-ci teste les
\emph{données} elles-mêmes, sur toute la population, pour y chercher ce
que les contrôles auraient dû empêcher. C'est l'outil le plus puissant de
l'auditeur public d'aujourd'hui~: un tableur suffit, et quelques
formules révèlent ce que des semaines de contrôle sur pièces ne
voient pas.

\separateur

% ============================================================================
\section{Le risque de fraude}
\label{sec:c12-risque}
% ============================================================================

\begin{boxdef}
\termecle{Fraude}\index{fraude}~: acte intentionnel destiné à obtenir un
avantage indu, par la tromperie ou la dissimulation. Elle se distingue de
l'erreur par l'intention, que l'auditeur ne peut pas établir seul.
\end{boxdef}

Le COSO demande à l'organisation d'évaluer le risque de fraude
(principe~8)~\cite{coso2013icfr}, et le guide commun du COSO et de l'ACFE
en détaille la démarche~\cite{coso2016fraude}. L'analyse classique de
Cressey réunit trois conditions, dites \termecle{triangle de la
fraude}\index{triangle de la fraude}~: une \emph{pression}, une
\emph{occasion}, une \emph{rationalisation}~\cite{cressey1953}.
L'auditeur ne peut agir que sur l'occasion~: c'est exactement ce que
créent les déficiences du contrôle interne.

\begin{tableaucours}{L{3.6cm}L{8.6cm}}{%
    \tableauHeader{Schéma} &
    \tableauHeader{Occasion dans le SI, et test de données}
}
Agent fictif & Un profil peut créer un agent et valider la paie~; tester
l'existence au fichier du personnel \\
Double paiement & Pas de blocage des doublons~; tester fournisseur,
montant, facture normalisée \\
Changement de bénéficiaire & RIB modifiable sans seconde validation~;
tester les RIB modifiés avant un paiement \\
Fractionnement & Seuil de mise en concurrence~; tester les montants juste
sous le seuil \\
Annulation de recette & Annulation sans second agent~; tester la
concentration des annulations \\
\end{tableaucours}
\vspace{-0.8em}
\begin{center}{\small\itshape Tableau~12.1. Schémas de fraude, occasions
et tests}\end{center}

\separateur

% ============================================================================
\section{Analyser un fichier de paie}
\label{sec:c12-paie}
% ============================================================================

\subsection{Trois questions}

\begin{enumerate}
    \item \textbf{Existence}~: chaque agent payé existe-t-il au fichier du
          personnel, en position d'activité~?
    \item \textbf{Unicité}~: chaque agent n'est-il payé qu'une fois, sur un
          compte qui n'est qu'à lui~?
    \item \textbf{Cohérence}~: l'âge, le grade, l'indice et les éléments de
          solde sont-ils compatibles entre eux et avec les règles~?
\end{enumerate}

\begin{boxoutils}{Tests de paie dans un tableur}
Feuille \texttt{PAIE}~: A matricule, B nom, C date de naissance, D compte
bancaire, E net. Feuille \texttt{PERSONNEL}~: A matricule.
\begin{lstlisting}
F2 (existence) : =SI(NB.SI(PERSONNEL!A:A;A2)=0;"ABSENT";"")
G2 (compte)    : =NB.SI(D:D;D2)
H2 (age)       : =DATEDIF(C2;DATE(2025;12;31);"y")
\end{lstlisting}
Un filtre sur \texttt{ABSENT}, sur \texttt{G~>~1} et sur \texttt{H~>~60}
donne les trois listes à examiner.
\end{boxoutils}

\begin{boxscenario}{Dossier CAPNO. Les trois listes}
\textbf{Constatation.} Sur les 4\,312 agents payés en décembre~2025~:
37 ne figurent pas au fichier du personnel (net cumulé
9\,842\,500~FCFA par mois)~; 12 comptes bancaires sont partagés par 27
agents~; 9 agents ont plus de 60 ans et sont toujours payés.

\medskip
\textbf{Analyse} (confiance moyenne). Les 37 absents sont des agents
fictifs \emph{probables}, pas certains~: un retard de mise à jour du
fichier du personnel peut en expliquer une partie. Les comptes partagés
peuvent être des comptes de couples. Chaque cas est examiné avec sa
pièce~; seuls les cas inexpliqués deviennent des constats.
\end{boxscenario}

\begin{boxalert}
\textbf{Violation typique.} Écrire au rapport «~37 agents fictifs~» à la
sortie du tableur.

\smallskip
\textbf{Règle~:} l'analyse de données produit des \emph{indices}. Le
constat vient après l'examen de chaque cas~: dossier individuel,
présence physique, acte de recrutement.
\end{boxalert}

\separateur

% ============================================================================
\section{Le fractionnement des commandes}
\label{sec:c12-fractionnement}
% ============================================================================

Quand un seuil déclenche une procédure plus lourde (mise en concurrence,
visa), la tentation est de découper une commande en plusieurs commandes
sous le seuil. Le test est simple~: compter les montants dans une bande
juste sous le seuil, par fournisseur et par période.

\begin{boxscenario}{Dossier MINÉQUIR. Juste sous le seuil}
\textbf{Constatation.} La DAF fixe à 5\,000\,000~FCFA le seuil de mise en
concurrence de ses bons de commande\footnote{Seuil interne fictif, propre
au cas~; il ne reproduit pas les seuils de la réglementation des marchés
publics.}. En 2025, 23 bons de commande sont compris entre 4\,800\,000 et
4\,999\,999~FCFA~; ils sont adressés à trois fournisseurs seulement.

\medskip
\textbf{Analyse} (confiance moyenne). La concentration dans une bande
étroite sous le seuil, sur peu de fournisseurs, est un indice classique
de fractionnement. Il faut examiner, pour chaque fournisseur, l'objet et
la date des bons~: des commandes de même objet à quelques jours
d'intervalle confirmeraient l'indice.
\end{boxscenario}

\separateur

% ============================================================================
\section{La loi de Benford}
\label{sec:c12-benford}
% ============================================================================

Dans de nombreuses séries de montants naturels, le premier chiffre n'est
pas uniformément réparti~: le 1 apparaît environ trois fois sur dix, le 9
moins d'une fois sur vingt~\cite{benford1938,nigrini2012benford}. La
proportion attendue du premier chiffre $d$ est $\log_{10}(1 + 1/d)$.

\begin{tableaucours}{L{3.0cm}ccccccccc}{%
    \tableauHeader{Premier chiffre} &
    \tableauHeader{1} & \tableauHeader{2} & \tableauHeader{3} &
    \tableauHeader{4} & \tableauHeader{5} & \tableauHeader{6} &
    \tableauHeader{7} & \tableauHeader{8} & \tableauHeader{9}
}
Attendu (\%) & 30,1 & 17,6 & 12,5 & 9,7 & 7,9 & 6,7 & 5,8 & 5,1 & 4,6 \\
\end{tableaucours}
\vspace{-0.8em}
\begin{center}{\small\itshape Tableau~12.2. Distribution de Benford du
premier chiffre}\end{center}

\begin{boxoutils}{Premier chiffre et proportion attendue}
\begin{lstlisting}
B2 (premier chiffre) : =GAUCHE(TEXTE(A2;"0");1)*1
Attendu pour d       : =LOG10(1+1/d)
Observe pour d       : =NB.SI(B:B;d)/NB(B:B)
\end{lstlisting}
\end{boxoutils}

\begin{boxinfo}
\textbf{Prudence.} La loi de Benford vaut pour des séries nombreuses,
étalées sur plusieurs ordres de grandeur, et non bornées par des seuils ou
des barèmes. Une paie (barème) ou des montants plafonnés ne la suivent
pas, sans aucune fraude. L'écart désigne où regarder~; il ne prouve rien.
\end{boxinfo}

\begin{boxscenario}{Dossier MINÉQUIR. Le chiffre 4}
\textbf{Constatation.} Parmi les montants des bons de commande de 2025, le
premier chiffre 4 apparaît dans 15,8~\% des cas, contre 9,7~\% attendus.

\medskip
\textbf{Analyse} (confiance faible à moyenne). L'excès de 4 est cohérent
avec les 23 bons compris entre 4,8 et 5~millions. Benford ne dit rien de
plus que le test de bande~; il le confirme par une autre voie.
\end{boxscenario}

\separateur

% ============================================================================
\section{De l'indice au constat}
\label{sec:c12-constat}
% ============================================================================

\begin{boxprocedure}
\textbf{Feuille d'analyse de données}~:
\begin{enumerate}
    \item fichier analysé (source, date, empreinte, totaux rapprochés)~;
    \item question posée et formule ou requête utilisée, recopiée
          intégralement~;
    \item résultat brut (nombre de cas, liste en annexe)~;
    \item examen de chaque cas, avec la pièce~: expliqué, inexpliqué, en
          attente~;
    \item constat sur les seuls cas inexpliqués~; \NEK{} sur les cas en
          attente.
\end{enumerate}
\end{boxprocedure}

\begin{boxmetier}{Vérificateur}
Gardez la formule exacte de chaque test dans le dossier. Lors du
contradictoire, l'entité refera votre calcul~; elle doit trouver le même
résultat.
\end{boxmetier}

\begin{boxmetier}{Juge des comptes}
Un paiement à un agent absent du fichier du personnel engage la
responsabilité de celui qui a payé sans vérifier l'existence du
créancier. Demandez quel contrôle d'application devait l'en empêcher.
\end{boxmetier}

\begin{boxmetier}{Auditeur interne}
Faites tourner chaque mois les trois tests de paie et le test de
fractionnement~: c'est un contrôle continu peu coûteux, qui relève du
pilotage (principe~16 du COSO).
\end{boxmetier}

\separateur

\begin{boxresume}
\begin{itemize}
    \item \textbf{La fraude suppose une intention}, que l'auditeur ne peut
          établir seul~; il agit sur l'occasion.
    \item \textbf{Chaque déficience ouvre un schéma}, et chaque schéma
          appelle un test de données.
    \item \textbf{Paie~:} existence, unicité, cohérence.
    \item \textbf{Fractionnement~:} une bande étroite sous le seuil, peu de
          fournisseurs.
    \item \textbf{Benford désigne, ne prouve pas}~; il ne s'applique pas à
          des séries bornées.
    \item \textbf{Un indice n'est pas un constat.} Chaque cas est examiné
          avec sa pièce.
\end{itemize}
\end{boxresume}

\begin{boxexo}[title={\ding{45}\quad Exercice~12.1. Le triangle \niveaubase}]
Pour le double paiement du MINÉQUIR, dites ce que l'auditeur peut
établir de chaque côté du triangle de la fraude, et ce qui ne relève pas
de lui.
\end{boxexo}

\begin{boxexo}[title={\ding{45}\quad Exercice~12.2. Benford \niveaubase}]
Calculez la proportion attendue du premier chiffre 1 et du premier
chiffre 9. Dites pourquoi on n'applique pas le test aux indemnités de
logement de la CAPNO, qui suivent un barème.
\end{boxexo}

\begin{boxexo}[title={\ding{45}\quad Exercice~12.3. Les comptes partagés \niveaumoyen}]
Proposez la démarche d'examen des 12 comptes partagés par 27 agents de la
CAPNO~: pièces, ordre d'examen, critères pour classer un cas comme
expliqué.
\end{boxexo}

\begin{boxexo}[title={\ding{45}\quad Exercice~12.4. Du tableur au rapport \niveauexpert}]
Après examen, 31 des 37 agents absents restent inexpliqués. Rédigez le
constat en cinq éléments, avec l'effet chiffré sur l'exercice 2025, et
dites ce que vous écrivez sur les 6 autres.
\end{boxexo}

\begin{boxinfo}
\textbf{Transition.} Il reste à regarder le système d'information de
plus haut~: qui le pilote, comment il est acheté, et à qui il est confié.
C'est l'objet du Chapitre~\ref{chap:gouvernance}.
\end{boxinfo}
